\documentclass{article}
\begin{document}
\section{Ouverture}
Le premier paragraphe commence ici. Un texte assez long pour remplir la page, phrase après phrase, sans rien de particulier. Un texte assez long pour remplir la page, phrase après phrase, sans rien de particulier. Un texte assez long pour remplir la page, phrase après phrase, sans rien de particulier. Un texte assez long pour remplir la page, phrase après phrase, sans rien de particulier. Un texte assez long pour remplir la page, phrase après phrase, sans rien de particulier. Un texte assez long pour remplir la page, phrase après phrase, sans rien de particulier. 

Le deuxième paragraphe suit. Un texte assez long pour remplir la page, phrase après phrase, sans rien de particulier. Un texte assez long pour remplir la page, phrase après phrase, sans rien de particulier. Un texte assez long pour remplir la page, phrase après phrase, sans rien de particulier. Un texte assez long pour remplir la page, phrase après phrase, sans rien de particulier. Un texte assez long pour remplir la page, phrase après phrase, sans rien de particulier. Un texte assez long pour remplir la page, phrase après phrase, sans rien de particulier. 
\subsection{Chapitre}
Un paragraphe venu d'un autre fichier, sur sa propre ligne du source,
et qui continue sur une deuxième ligne du source, assez longue pour passer à la ligne.

\newpage
\section{Suite}
Le texte de la seconde page est court.
Une formule $a^2 + b^2 = c^2$ au milieu.
\end{document}
